\documentclass[10pt]{article}
\usepackage[T1]{fontenc}
\usepackage[utf8]{inputenc}
\usepackage[french]{babel}
\usepackage{amsmath,amssymb,mathrsfs,mathpazo}
\usepackage{geometry}
\geometry{a4paper,margin=2cm}
\usepackage{graphicx,array,tabularx,tabularray,colortbl,multirow,multicol,booktabs,enumitem,pifont,ulem,fancybox,diagbox,cancel}
\usepackage{pstricks,pstricks-add,pst-plot,pst-tree,pst-node,pst-text,pst-func,pst-math,pst-bspline,pst-3dplot,pst-eucl,pst-solides3d}
\let\pscalculate\undefined
\usepackage{pst-bezier}
\usepackage{tikz,pgfplots,tkz-tab}
\usetikzlibrary{arrows,arrows.meta,patterns,calc,positioning,trees,decorations.markings,shapes.geometric,shapes.misc,shapes.arrows,intersections,decorations.pathmorphing,decorations.pathreplacing,backgrounds,shadows,babel}
\pgfplotsset{compat=1.18}
\usepgfplotslibrary{fillbetween}
\usepackage[np]{numprint}
\usepackage{listings,makecell,pas-tableur,twoopt}
\let\barre\undefined
\usepackage[pstricks]{tablvar}
\definecolor{darkgreen}{RGB}{0,100,0}
\providecommand{\red}{\color{red}}
\providecommand{\blue}{\color{blue}}
\providecommand{\green}{\color{green}}
\providecommand{\black}{\color{black}}
\providecommand{\white}{\color{white}}
\providecommand{\gray}{\color{gray}}
\providecommand{\cyan}{\color{cyan}}
\providecommand{\magenta}{\color{magenta}}
\providecommand{\yellow}{\color{yellow}}
\providecommand{\violet}{\color{violet}}
\providecommand{\orange}{\color{orange}}
\providecommand{\euro}{\texteuro}
\providecommand{\EUR}{\texteuro}
\providecommand{\ieme}{\textsuperscript{e}}
\providecommand{\degres}{\ensuremath{^\circ}}
\providecommand{\degree}{\ensuremath{^\circ}}
\providecommand{\vv}[1]{\overrightarrow{#1}}
\providecommand{\Coord}[1]{#1}
\providecommand{\decosix}{}
\providecommand{\starredbullet}{\ensuremath{\bullet}}
\pagestyle{empty}
\setlength{\parindent}{0pt}
\usepackage{adjustbox}\DeclareUnicodeCharacter{2236}{:}\DeclareUnicodeCharacter{2190}{\ensuremath{\leftarrow}}
\providecommand{\R}{\mathbb{R}}
\providecommand{\N}{\mathbb{N}}
\providecommand{\Z}{\mathbb{Z}}
\providecommand{\Q}{\mathbb{Q}}
\providecommand{\D}{\mathbb{D}}
\providecommand{\C}{\mathbb{C}}
\providecommand{\e}{{\mathrm{e}}}
\providecommand{\vect}[1]{\overrightarrow{#1}}
\providecommand{\vectt}[1]{\overrightarrow{\mathrm{#1}}}
\providecommand{\vvt}[1]{\overrightarrow{\mathrm{#1}}}
\providecommand{\np}[2][]{{#2~#1}}
\providecommand{\vv}[1]{{\overrightarrow{#1}}}
\renewcommand{\d}{{\mathrm{d}}}
\providecommand{\degres}{{^\circ}}
\providecommand{\degre}{{^\circ}}
\providecommand{\blue}{\color{blue}}
\providecommand{\red}{\color{red}}
\providecommand{\enskip}{\quad}
\providecommand{\llbracket}{\lbrack\!\lbrack}
\providecommand{\rrbracket}{\rbrack\!\rbrack}
\providecommand{\euro}{{\text{€}}}
\providecommand{\no}{{\text{n°}}}
\providecommand{\interval}[2]{{[#1;#2]}}
\providecommand{\intervFO}[2]{{[#1;#2[}}
\providecommand{\dotfill}{\ldots}

\providecommand{\ds}{\displaystyle}
\providecommand{\textding}[1]{#1}
\providecommand{\Oij}{\(\mathrm{O};\vec{\imath},\vec{\jmath}\)}
\providecommand{\Oijk}{\(\mathrm{O};\vec{\imath},\vec{\jmath},\vec{k}\)}
\providecommand{\Ouv}{\(\mathrm{O};\vec{u},\vec{v}\)}
\providecommand{\intervFF}[2]{\ensuremath{[#1;#2]}}
\renewcommand{\euro}{\ifmmode\text{\texteuro}\else\texteuro\fi}
\DeclareUnicodeCharacter{00B0}{\ifmmode^{\circ}\else\textdegree\fi}
\DeclareUnicodeCharacter{20AC}{\ifmmode\text{\texteuro}\else\texteuro\fi}
% Énoncés : Éducation nationale. Transcription et corrigés : APMEP, Denis Vergès et contributeurs. Découpage et classement : CoopMaths, Sébastien Lozano, Éric Elter et contributeurs. Adaptation de la présentation : Maths Explorer.
% CC BY-SA : https://coopmaths.fr/www/about/
% https://coopmaths.fr/alea/static/bac/2021/tex/bac_2021_06_polynesie_3_cor.tex
\begin{document}
\noindent\textit{APMEP / CoopMaths. CC BY-SA. Présentation adaptée par Maths Explorer.}\par\medskip
\textbf{Commun à tous les candidats}







\begin{center}
\psset{unit=1cm}
\begin{pspicture}(-4,-2.75)(4.5,5.1)
\psline[linewidth=1.25pt]{->}(0,5.)\uput[u](0,5){$z$}
\psline[linewidth=1.25pt]{->}(4.5,-1.8)\uput[dr](4.5,-1.8){$y$}
\psline[linewidth=1.25pt]{->}(-4,-2)\uput[dl](-4,-2){$x$}
\pspolygon[linewidth=1.25pt](-2.4,-1.16)(0.5,-2.4)(2.85,-1.1)(2.85,2.6)(0,3.6)(-2.4,2.3)
\psline[linewidth=1.25pt](-2.4,2.3)(0.5,1.06)(2.85,2.6)
\psline[linewidth=1.25pt](0.5,1.06)(0.5,-2.4)
\psline[linestyle=dashed,linewidth=1.25pt](-2.4,-1.16)(2.85,2.6)
\psdots(-0.65,0.093) \uput[ul](-0.65,0.093){M}
\uput[d](-2.4,-1.16){B} \uput[d](0.5,-2.4){C} \uput[d](2.85,-1.1){D} 
\uput[r](2.85,2.6){H} \uput[ur](0,3.6){E} \uput[l](-2.4,2.3){F} 
\uput[d](0,0){A} \uput[u](0.5,1.06){G}
\end{pspicture}
\end{center}


\begin{enumerate}
\item 
Dans le repère $\left(\text{A}~;~\vect{\text{AB}},~\vect{\text{AD}},~\vect{\text{AE}}\right)$, on a 

B(1~;~0~;~0), \, D(0~;~1~;~0), \, E(0~;~0~;~1),\, G(1~;~1~;~1), \, H(0~;~1~;~1).
\item 
	\begin{enumerate}
		\item{}  
[EG], [GD] et [ED] sont les hypoténuses de triangles rectangles isocèles de côté 1, donc EG = GD = ED $= \sqrt{2}$ : le triangle EGD est équilatéral.
		\item 
		

Puisque $c = \sqrt{2}$, on a $\mathcal{A}(\text{EGD}) = \dfrac{\sqrt{3}}{4}\times \left(\sqrt{2}\right)^2 = \dfrac{\sqrt{3}}{4}\times 2 = \dfrac{\sqrt{3}}{2}$.
	\end{enumerate}
	
\item 

On a $\vect{\text{BH}}\begin{pmatrix}- 1\\1\\1\end{pmatrix}$, donc $\vect{\text{BM}} = \dfrac{1}{3}\vect{\text{BH}}\begin{pmatrix}- \frac{1}{3}\\[3pt]\frac{1}{3}\\[3pt]\frac{1}{3}\end{pmatrix} = \begin{pmatrix}x_{\text{M}} - 1\\y_{\text{M}}- 0\\z_{\text{M}}- 0\end{pmatrix}$.

On a donc : 
$x_{\text{M}} = 1 - \dfrac{1}{3} = \dfrac{2}{3}$ ;\,
$y_{\text{M}} = \dfrac{1}{3} $ ; \,
$z_{\text{M}} = \dfrac{1}{3}$.

Conclusion : M a pour coordonnées $\left(\dfrac{2}{3}~;~\dfrac{1}{3}~; ~\dfrac{1}{3}\right)$.
\item
	\begin{enumerate}
		\item 
		
On a $\vect{\text{EG}}\begin{pmatrix}1\\1\\0\end{pmatrix}$ : donc $\vect{n} \cdot \vect{\text{EG}} = - 1 +1 + 0 = 0$ :

$\vect{\text{ED}}\begin{pmatrix}0\\1\\-1\end{pmatrix}$ : donc $\vect{n} \cdot \vect{\text{ED}} =  0 + 1 - 1  = 0$.

Conclusion : $\vect{n}$ est orthogonal à deux vecteurs non colinéaires du plan (EGD) : c'est un vecteur normal à ce plan.
		\item 
		On sait qu'un plan d'équation $ax + by + cz + d = 0$ a un vecteur normal de coordonnées $(a~;~b~;~c)$, donc :
		
$P(x~;~y~;~z) \in (\text{EGD})  \iff -x + y + z + d = 0$.

Comme E(0~;~0~;~1)$\in (\text{EGD}) \iff 0 + 0 + 1 + d = 0 \iff d = - 1$.

Finalement : le plan (EGD) a pour équation $-x + y + z - 1 = 0$.
		\item 
		









La droite $\mathcal{D}$ contient M et a pour vecteur directeur $\vect{n}$ vecteur normal au plan (EGD), donc avec $\vect{\text{M}P}\begin{pmatrix}x - \frac{2}{3}\\[3pt]y - \frac{2}{3}\\[3pt]z - \frac{1}{3}\end{pmatrix}$:

$P(x~;~y~;~z) \in \mathcal{D} \iff \vect{\text{M}P} = t \vect{n}$, avec $t \in \R$, soit :

$\left\{\begin{array}{l c l}
x - \frac{2}{3}&=& - t\\[3pt]
y - \frac{1}{3}&=&t\\[3pt]
z - \frac{1}{3}&=&t
\end{array}\right., \, t \in \R \iff \left\{\begin{array}{l c l}
x	&=&\frac{2}{3} - t\\[3pt]
y	&=&\frac{1}{3} + t\\[3pt]
z 	&=&\frac{1}{3} + t
\end{array}\right., \, t \in \R$. 
	\end{enumerate}
\item Le cube ABCDEFGH est représenté ci-dessus selon une vue qui permet de mieux percevoir la pyramide GEDM, en gris sur la figure :

\begin{center}
\psset{unit=1cm}
\begin{pspicture}(-7,-0.5)(2,6.8)
\psline[linewidth=2pt,fillstyle=solid,fillcolor=lightgray](-4.2,0.6)(-0.5,2.6)(0,4.2)(-2.9,6.2)
\pspolygon[linestyle=dotted,linewidth=2pt,fillstyle=solid,fillcolor=gray](-2.9,6.2)(-4.2,0.6)(-0.5,2.6)
\psline[linewidth=2pt](-4.2,0.6)(0,4.2)(-2.9,6.2)
\psline[linewidth=1.25pt]{->}(1.9,2.3)\uput[ur](1.9,2.3){$x$}
\psline[linewidth=1.25pt]{->}(-6.4,0.9)\uput[ul](-6.4,0.9){$y$}
\psline[linewidth=1.25pt]{->}(0,6.2)\uput[u](0,6.2){$z$}
\pspolygon[linewidth=1.25pt](0,0)(1.3,1.55)(1.3,5.8)(-2.9,6.2)(-4.2,4.6)(-4.2,0.6)
\psline[linewidth=1.25pt](-4.2,4.6)(0,4.2)(1.3,5.8)

\psline[linestyle=dashed,linewidth=1.25pt](1.3,1.55)(-2.9,2.1)(-2.9,6.2)
\psline[linestyle=dashed,linewidth=1.25pt](-2.9,2.1)(-4.2,0.6)

\psline[linestyle=dotted,linewidth=2pt](-0.5,2.6)(0,4.2)

\uput[d](-0.5,2.6){M}
\uput[r](1.3,1.55){B} \uput[ul](-2.9,2.1){C} \uput[d](-4.2,0.6){D} 
\uput[l](-4.2,4.6){H} \uput[r](0,4.2){E} \uput[r](1.3,5.8){F} 
\uput[d](0,0){A} \uput[u](-2.9,6.2){G}
\end{pspicture}
\end{center}



	\begin{enumerate}
		\item 
		

On a vu que la droite $\mathcal{D}$ contient M et est perpendiculaire au plan (EBD) : c'est donc la hauteur de la pyramide GEDM issue de M. Le pied de cette hauteur K appartient donc à $\mathcal{D}$ et au plan (EGD) ; ses coordonnées vérifient donc les équations :

$\left\{\begin{array}{l c l}
x	&=&\frac{2}{3} - t\\[3pt]
y	&=&\frac{1}{3} + t\\[3pt]
z 	&=&\frac{1}{3} + t\\[3pt]
- x + y + z - 1&=&0
\end{array}\right.\Rightarrow -\left(\frac{2}{3} - t \right) + \frac{1}{3} + t + \frac{1}{3} + t - 1 = 0 \iff 3t = 1 \iff t = \frac{1}{3}$.

En remplaçant $t$ par $\frac{1}{3}$ dans l'équation paramétrique de $\mathcal{D}$, on obtient :

$\left\{\begin{array}{l c l}
x	&=&\frac{2}{3} - \frac{1}{3} = \frac{1}{3}\\[3pt]
y	&=&\frac{1}{3} + \frac{1}{3} = \frac{2}{3}\\[3pt]
z 	&=&\frac{1}{3} + \frac{1}{3} = \frac{2}{3}
\end{array}\right.$

Conclusion : K a pour coordonnées $\left(\dfrac{1}{3}~;~\dfrac{2}{3}~;~\dfrac{2}{3} \right)$.
		\item 
On en déduit $\vect{\text{KM}}\begin{pmatrix}\frac{1}{3}\\[3pt]- \frac{1}{3}\\[3pt]- \frac{1}{3}\end{pmatrix}$.

D'où KM$^2 = \dfrac{1}{9} + \dfrac{1}{9} + \dfrac{1}{9} = \dfrac{3}{9}$ ; on en déduit KM$ = \dfrac{\sqrt{3}}{3}$.
		




Comme $\mathcal{A}(\text{EGD}) = \frac{\sqrt{3}}{2}$, le volume de la pyramide GEDM est  : $\dfrac{\frac{\sqrt{3}}{2}\times \frac{\sqrt{3}}{3}}{3} = \dfrac{1}{6}$.
	\end{enumerate}
\end{enumerate}

\bigskip

\end{document}
